\section{El error de torres en los cuatro vértices}
Sea $\mathcal P(X)=\log\per(D_{-X})$, con $X$ un vector indicador de las marcas impuestas.
Se trata del menor del \emph{exterior} $D$: los otros sitios marcados siguen ausentes.
Por los apartados 2 y 5, con $m=j-|X|$,
\[
 \mathcal P(X)=R(X)+A(X)+\tfrac12\log(4\pi j)+u_T(D_{-X})+R_T(D_{-X}) .
\]
En efecto $t\log t=0$ si $t=0,1$, y el cambio del modo nulo aporta
$\tfrac12\log(m/j)$ al límite de $A$. Se usa aquí el error puro de torres
$R_T=\log\per-L_T-u_T$. R9, apartados 6--8, da una variación $O(s^3)$ de $R_T$
al retirar una fila y una columna, también en los sucesivos menores de tamaño fijo.
La estabilidad de $F_2$ de R9, apartado 8, sólo usa la restricción de los índices y la
corrección de márgenes; el cruce suprimido puede estar ausente. Así una diferencia
de vértices de $R_T$ cuesta $O(s^3)$, y lo mismo una diferencia doble.

Para el corrector $u_T$ hace falta algo más que su variación a un paso.
Normalizamos ahora las medidas tras retirar hasta dos grupos de dos índices en cada lado:
\[
 \mu_h=\frac{\mu^0-h_1\sum_{i\in u}\delta_i-h_2\sum_{i\in v_*}\delta_i}
                  {1-2h_1-2h_2},\qquad s_m=\frac{s}{1-2h_1-2h_2}.
\]
Se entiende un índice de fila o columna por marca en su lado correspondiente.
En $h_g=0,s$ obtenemos exactamente las medidas uniformes de los cuatro menores.
Con el mismo núcleo basal resolvemos $v(1-s_m v)=z\exp(\alpha+\beta)$ con ambos márgenes
uno en $\mu_h$. La prueba del apartado 3 da $v$ acotado y holomorfo en un bidisco fijo.
En los vértices $P=s_m v$ al restringir a las filas y columnas supervivientes.
La fórmula de R9, apartados 3 y 7, da
$u_T=s_m[-1/12+\phi(v)]$. Detallamos por qué $\phi$ permanece holomorfa y acotada.
Para $b=v-1$ intervienen $G=bb^t$ o $b^tb$, $T=(I-G)^{-1}G$,
$V=(I-G)^{-1}b$, $U=V-b$ con el lado correspondiente.
Las transpuestas y convoluciones se toman bilinealmente en las medidas.
En el origen ($s=h=0$) la brecha acota el inverso en $L^2$ y la identidad
$(I-G)^{-1}=I+G+G(I-G)^{-1}G$ lo acota en supremo por las cotas del núcleo.
La variación total de la medida y el cambio de $v$ son pequeños;
en supremo la inversa sigue acotada por Neumann en discos menores.
Aplicarla a cada columna de $G$ o $b$ acota el núcleo de $T,V$ incluida su diagonal.
En R9 $\phi$ es suma de trazas de productos suavizados, integrales de cubos por entradas,
y formas cuadráticas con la diagonal de $T-G$. Cada una integra un número fijo de
núcleos así acotados, por lo que la fórmula conserva holomorfía y cota uniforme.
En un índice de masa cero sólo se extiende el núcleo; sus integrales son las del menor.
La cota $|u_T|\le Cs$ en el bidisco da por Cauchy las mismas cotas para las primeras
y segundas derivadas en masas en un disco menor. El paso de masa $s$ cuesta $O(s^2)$
y la diferencia mixta de los dos grupos cuesta $O(s^3)$. La primera afirmación también
vale para cualquier subconjunto fijo, ajustando el número de índices por grupo.

Sea $r=\nabla R(0)$ y $g(X)=\mathcal P(X)-\mathcal P(0)-r\cdot X$.
Escribamos $H=R''(0)$, $U_3=R'''(0)$, $a=A'(0)$ y $h=H[u,v_*]$.
En todos los vértices, por el apartado 6 y las primeras diferencias de torres,
\[
 g(X)=a[X]+\tfrac12 H[X,X]+O(s^2).
\]
Para los dos grupos, usando ahora su diferencia doble, resulta
\[
 \Delta_{u,v_*}g=h+A''(0)[u,v_*]
       +\tfrac12(U_3[u,u,v_*]+U_3[u,v_*,v_*])+O(s^3).
\]
Ésta es la precisión que se pierde si se usa sólo una cota sobre cada valor de $R_T+u_T$.
